% ---
% school: "college"
% level: "8eme"
% section: null
% chapter: "العمليات والحساب على الأعداد الكسرية"
% source: "Axiomaths"
% title: "مجموع تسعة أعداد كسرية"
% country: "Tunisie"
% year: 2026
% professor: "Axiomaths"
% difficulty: "Difficile"
% points: null
% tags: ["fractions", "somme-telescopique", "decomposition"]
% ---
% Exercice 14/18 de la série 8ème « العمليات والحساب على الأعداد الكسرية » (statement + \axsol + solution). Compiler avec XeLaTeX + axiomath.sty (+ axiomaths-logo.png).
\documentclass[12pt,a4paper]{article}
\usepackage[level=8b,color=15803D]{axiomath}
\axsetup{title={العمليات والحساب على الأعداد الكسرية}, corrige}

\begin{document}

\begin{corr}[title={مجموع تسعة أعداد كسرية}, stars=4, source={Axiomaths}]
نعتبر المجموع :
\[S=\dfrac{1}{2}+\dfrac{1}{6}+\dfrac{1}{12}+\dfrac{1}{20}+\dfrac{1}{30}+\dfrac{1}{42}+\dfrac{1}{56}+\dfrac{1}{72}+\dfrac{1}{90}\]
\begin{enumerate}
  \item تحقق أن $\dfrac{1}{2}=1-\dfrac{1}{2}$ ، $\dfrac{1}{6}=\dfrac{1}{2}-\dfrac{1}{3}$ ، $\dfrac{1}{12}=\dfrac{1}{3}-\dfrac{1}{4}$ و $\dfrac{1}{20}=\dfrac{1}{4}-\dfrac{1}{5}$.
  \item نقبل أن $\dfrac{1}{n\times(n+1)}=\dfrac{1}{n}-\dfrac{1}{n+1}$ لكل عدد صحيح طبيعي $n$ غير معدوم.
  اكتب كل حد من $S$ على شكل فرق عددين كسريين. [لاحظ أن $6=2\times3$ و $12=3\times4$]
  \item استنتج أن $S=\dfrac{9}{10}$.
  \item احسب $T=\dfrac{1}{6}+\dfrac{1}{12}+\dfrac{1}{20}+\dfrac{1}{30}+\dfrac{1}{42}+\dfrac{1}{56}+\dfrac{1}{72}+\dfrac{1}{90}$.
  \item احسب المجموع $\dfrac{1}{1\times2}+\dfrac{1}{2\times3}+\dfrac{1}{3\times4}+\cdots+\dfrac{1}{99\times100}$.
\end{enumerate}
\axsol
\begin{enumerate}[itemsep=14pt]
  \item
  \[1-\dfrac{1}{2}=\dfrac{2}{2}-\dfrac{1}{2}=\dfrac{1}{2}\qquad\dfrac{1}{2}-\dfrac{1}{3}=\dfrac{3}{6}-\dfrac{2}{6}=\dfrac{1}{6}\]
  \[\dfrac{1}{3}-\dfrac{1}{4}=\dfrac{4}{12}-\dfrac{3}{12}=\dfrac{1}{12}\qquad\dfrac{1}{4}-\dfrac{1}{5}=\dfrac{5}{20}-\dfrac{4}{20}=\dfrac{1}{20}\]
  \item المقامات هي $1\times2$ ، $2\times3$ ، $3\times4$ ، $\dots$ ، $9\times10$ :
  \[\dfrac{1}{2}=1-\dfrac{1}{2}\qquad\dfrac{1}{6}=\dfrac{1}{2}-\dfrac{1}{3}\qquad\dfrac{1}{12}=\dfrac{1}{3}-\dfrac{1}{4}\]
  \[\dfrac{1}{20}=\dfrac{1}{4}-\dfrac{1}{5}\qquad\dfrac{1}{30}=\dfrac{1}{5}-\dfrac{1}{6}\qquad\dfrac{1}{42}=\dfrac{1}{6}-\dfrac{1}{7}\]
  \[\dfrac{1}{56}=\dfrac{1}{7}-\dfrac{1}{8}\qquad\dfrac{1}{72}=\dfrac{1}{8}-\dfrac{1}{9}\qquad\dfrac{1}{90}=\dfrac{1}{9}-\dfrac{1}{10}\]
  \item نعوّض كل حد بفرقه :
  \[S=\left(1-\dfrac{1}{2}\right)+\left(\dfrac{1}{2}-\dfrac{1}{3}\right)+\left(\dfrac{1}{3}-\dfrac{1}{4}\right)+\cdots+\left(\dfrac{1}{9}-\dfrac{1}{10}\right)\]
  \[S=1-\dfrac{1}{2}+\dfrac{1}{2}-\dfrac{1}{3}+\dfrac{1}{3}-\dfrac{1}{4}+\cdots+\dfrac{1}{9}-\dfrac{1}{10}\]
  كل عدد كسري يظهر مرتين بإشارتين متعاكستين فيُحذف، ولا يبقى إلا الحد الأول والحد الأخير :
  \[S=1-\dfrac{1}{10}=\dfrac{9}{10}\]
  \item $T$ هو $S$ بعد حذف الحد الأول $\dfrac{1}{2}$ :
  \[T=S-\dfrac{1}{2}=\dfrac{9}{10}-\dfrac{5}{10}=\dfrac{4}{10}=\dfrac{2}{5}\]
  \item نفس الطريقة : المجموع هو $1-\dfrac{1}{100}$ :
  \[\dfrac{1}{1\times2}+\cdots+\dfrac{1}{99\times100}=1-\dfrac{1}{100}=\dfrac{99}{100}\]
\end{enumerate}
\textbf{الفكرة :} كتابة كل حد على شكل فرق تجعل الحدود تُحذف مثنى مثنى (مجموع تلسكوبي).

\textbf{انتبه :} لا ننسى الحد الأول $1$ ولا الحد الأخير $-\dfrac{1}{10}$ اللذين لا يُحذفان.
\end{corr}

\end{document}
